\usepackage{amsmath,amssymb,amsthm,mathtools,bm}
\usepackage[a4paper,left=2.7cm,right=2.7cm,top=2.8cm,bottom=2.8cm]{geometry}
\usepackage{siunitx}
\usepackage{booktabs,longtable,array}
\usepackage{graphicx}
\usepackage{tikz,pgfplots}
\usepackage{xcolor}
\usepackage[unicode,pdfencoding=auto,psdextra]{hyperref}
\usepackage{bookmark}
\usepackage{cleveref}
\usepackage{enumitem}
\usepackage[most]{tcolorbox}
\usepackage{tabularx}

\pgfplotsset{compat=1.18}
\usetikzlibrary{arrows.meta,positioning,calc,fit,shapes.geometric}

\hypersetup{
  colorlinks=true,
  linkcolor=blue!60!black,
  citecolor=green!40!black,
  urlcolor=magenta!60!black,
  pdftitle={PCSEL 理论、仿真与器件物理导论},
  pdfauthor={Feiyang Wu}
}

\pdfstringdefDisableCommands{
  \def\bm#1{#1}
  \def\mathbf#1{#1}
  \def\mathsf#1{#1}
  \def\mathrm#1{#1}
  \def\SI#1#2{#1 #2}
  \def\si#1{#1}
  \def\texorpdfstring#1#2{#2}
  \def\Gamma{Gamma}
  \def\Omega{Omega}
  \def\omega{omega}
  \def\lambda{lambda}
  \def\alpha{alpha}
  \def\beta{beta}
  \def\mu{mu}
  \def\varepsilon{epsilon}
  \def\to{->}
  \def\quad{ }
  \def\qquad{ }
  \def\hspace#1{ }
  \def\path#1{#1}
}

\setlist[itemize]{leftmargin=2em}
\setlist[enumerate]{leftmargin=2.5em}
\setcounter{secnumdepth}{3}
\setcounter{tocdepth}{2}
\newcommand{\basicq}{\textbf{[基础]}}
\newcommand{\midq}{\textbf{[进阶]}}
\newcommand{\hardq}{\textbf{[综合]}}

\providecommand{\crefpairconjunction}{和}
\providecommand{\crefmiddleconjunction}{、}
\providecommand{\creflastconjunction}{和}
\providecommand{\crefrangeconjunction}{至}

\crefformat{chapter}{第#2#1#3章}
\Crefformat{chapter}{第#2#1#3章}
\crefrangeformat{chapter}{第#3#1#4章至第#5#2#6章}
\Crefrangeformat{chapter}{第#3#1#4章至第#5#2#6章}
\crefmultiformat{chapter}{第#2#1#3章}{、第#2#1#3章}{、第#2#1#3章}{和第#2#1#3章}
\Crefmultiformat{chapter}{第#2#1#3章}{、第#2#1#3章}{、第#2#1#3章}{和第#2#1#3章}

\crefformat{section}{第#2#1#3节}
\Crefformat{section}{第#2#1#3节}
\crefrangeformat{section}{第#3#1#4节至第#5#2#6节}
\Crefrangeformat{section}{第#3#1#4节至第#5#2#6节}
\crefmultiformat{section}{第#2#1#3节}{、第#2#1#3节}{、第#2#1#3节}{和第#2#1#3节}
\Crefmultiformat{section}{第#2#1#3节}{、第#2#1#3节}{、第#2#1#3节}{和第#2#1#3节}

\crefformat{subsection}{第#2#1#3小节}
\Crefformat{subsection}{第#2#1#3小节}
\crefrangeformat{subsection}{第#3#1#4小节至第#5#2#6小节}
\Crefrangeformat{subsection}{第#3#1#4小节至第#5#2#6小节}
\crefmultiformat{subsection}{第#2#1#3小节}{、第#2#1#3小节}{、第#2#1#3小节}{和第#2#1#3小节}
\Crefmultiformat{subsection}{第#2#1#3小节}{、第#2#1#3小节}{、第#2#1#3小节}{和第#2#1#3小节}

\crefformat{equation}{式~#2(#1)#3}
\Crefformat{equation}{式~#2(#1)#3}
\crefrangeformat{equation}{式~#3(#1)#4至式~#5(#2)#6}
\Crefrangeformat{equation}{式~#3(#1)#4至式~#5(#2)#6}
\crefmultiformat{equation}{式~#2(#1)#3}{和式~#2(#1)#3}{、式~#2(#1)#3}{和式~#2(#1)#3}
\Crefmultiformat{equation}{式~#2(#1)#3}{和式~#2(#1)#3}{、式~#2(#1)#3}{和式~#2(#1)#3}

\crefformat{figure}{图#2#1#3}
\Crefformat{figure}{图#2#1#3}
\crefrangeformat{figure}{图#3#1#4至图#5#2#6}
\Crefrangeformat{figure}{图#3#1#4至图#5#2#6}
\crefmultiformat{figure}{图#2#1#3}{和图#2#1#3}{、图#2#1#3}{和图#2#1#3}
\Crefmultiformat{figure}{图#2#1#3}{和图#2#1#3}{、图#2#1#3}{和图#2#1#3}

\crefformat{table}{表#2#1#3}
\Crefformat{table}{表#2#1#3}
\crefrangeformat{table}{表#3#1#4至表#5#2#6}
\Crefrangeformat{table}{表#3#1#4至表#5#2#6}
\crefmultiformat{table}{表#2#1#3}{和表#2#1#3}{、表#2#1#3}{和表#2#1#3}
\Crefmultiformat{table}{表#2#1#3}{和表#2#1#3}{、表#2#1#3}{和表#2#1#3}

\newcolumntype{Y}{>{\raggedright\arraybackslash}X}

\newtheorem{theorem}{定理}[chapter]
\newtheorem{proposition}[theorem]{命题}
\newtheorem{lemma}[theorem]{引理}
\newtheorem{corollary}[theorem]{推论}
\theoremstyle{definition}
\newtheorem{definition}[theorem]{定义}
\newtheorem{remark}[theorem]{注记}
\newtheorem{example}[theorem]{例}

\tcbset{
  colback=white,
  colframe=black!60,
  boxrule=0.6pt,
  arc=1pt
}
\newtcolorbox{IntuitionBox}{title=直觉框,fonttitle=\heiti,colframe=blue!60!black}
\newtcolorbox{RigorousBox}{title=严格表述,fonttitle=\heiti,colframe=red!60!black}
\newtcolorbox{ExampleBox}{title=例子,fonttitle=\heiti,colframe=green!40!black}
\newtcolorbox{PitfallBox}{title=易错点,fonttitle=\heiti,colframe=orange!80!black}

\newcommand{\kvec}{\bm{k}}
\newcommand{\gvec}{\bm{G}}
\newcommand{\rvec}{\bm{r}}
\newcommand{\E}{\bm{E}}
\newcommand{\Hf}{\bm{H}}
\newcommand{\D}{\bm{D}}
\newcommand{\B}{\bm{B}}
\newcommand{\J}{\bm{J}}
\newcommand{\p}{\partial}
\newcommand{\ii}{\mathrm{i}}
\newcommand{\ee}{\mathrm{e}}
\newcommand{\dd}{\mathrm{d}}
\newcommand{\avg}[1]{\left\langle #1 \right\rangle}
\newcommand{\TODO}[1]{\textcolor{red}{TODO: #1}}
\numberwithin{equation}{chapter}

% Keep the build log focused on substantive problems rather than harmless
% CJK font-mapping artifacts and narrow-box diagnostics from dense tables.
\tracinglostchars=0
\hbadness=10000
\vbadness=10000
